% Generated draft: requires mathematical and source review.
\documentclass[11pt]{article}
\usepackage[margin=25mm]{geometry}
\usepackage{fontspec}
\setmainfont{Latin Modern Roman}
\usepackage{amsmath,amssymb,mathtools,unicode-math}
\setmathfont{Latin Modern Math}
\usepackage{enumitem,needspace,xurl,hyperref}
\setlength{\parindent}{0pt}
\setlength{\parskip}{5pt}
\setlength{\emergencystretch}{3em}
\newcommand{\N}{\mathbb{N}}
\newcommand{\Z}{\mathbb{Z}}
\newcommand{\Q}{\mathbb{Q}}
\newcommand{\R}{\mathbb{R}}
\newcommand{\C}{\mathbb{C}}
\newcommand{\F}{\mathbb{F}}
\newcommand{\A}{\mathbb{A}}
\newcommand{\PP}{\mathbb{P}}
\newcommand{\E}{\mathbb{E}}
\newcommand{\eps}{\varepsilon}
\newcommand{\dd}{\,\mathrm{d}}
\begin{document}
\section*{UnsolvedMath ID: 30001963}
\textbf{Source problem code:} catalogue entry
\subsection{Definitions and mathematical statement}
The supplied record concerns tacit mathematical skills, arts, or knowledge: those not specified or explicated in mathematical sources. This description does not supply a mathematical definition of tacit knowledge or of its recovery.

The original question is: \emph{How can historians recover skills, arts, or knowledge that were not specified or explicated in mathematical sources?}

The supplied statement-verification note characterizes this as a historical, pedagogical, or methodological discussion prompt, rather than an unresolved mathematical proposition. It assigns the statement status \texttt{unrecoverable} and supplies no separate clean statement. No class of historical sources, admissible evidence, recovery procedure, criterion of successful recovery, or mathematical quantifiers is specified. Consequently, a precise mathematical proposition cannot be extracted without adding substantive assumptions; the original methodological question is retained for provenance.
\subsection{Short English statement}
The question asks how historians can reconstruct mathematical skills, arts, or knowledge left unexpressed in historical mathematical sources. The supplied record treats it as a methodological discussion question; a mathematical formalization would require additional definitions and a specified target.
\subsection{Sources}
\begin{itemize}
\item Local UnsolvedMath registry, supplied record with permanent ID \(30001963\), exact source problem code \texttt{catalogue entry}, titled \emph{Recovering Tacit Mathematical Skills from Historical Sources}. The record associates the question with 2012 and names no individual proposer.
\item \emph{Explicit Versus Tacit Knowledge in Mathematics} (2012), Oberwolfach Reports, pp.~131--210, DOI 10.4171/owr/2012/04. Supplied source citation and literature-source link, labeled ``original Oberwolfach report'': \url{https://doi.org/10.4171/owr/2012/04}. The supplied background describes Oberwolfach Reports as documenting mathematical workshops and their open-problem discussions.
\item Literature assessment dated \texttt{2026-08-21}: the supplied curation labels the record \texttt{open}, describes it as a historiographical question about recovering tacit mathematical skills rather than a resolved mathematical theorem, and reports that no definitive answer was located. This is supplied dated curation, not an independently verified assessment. The supplied background requests independent verification of time-sensitive status and cited claims.
\end{itemize}
\end{document}
